\section*{Introduction}
Le marché des véhicules autonomes connaissant un essor considérable, les avancées technologiques majeures des dernières décennies ont permis de concrétiser ce qui relevait autrefois de la science-fiction. En associant des capteurs de perception de pointe, des architectures embarquées capables de traiter une importante masse de données en temps réel et, plus récemment, des algorithmes d'intelligence artificielle, les véhicules modernes sont désormais en mesure de naviguer de manière autonome au sein de scénarios de circulation complexes.\\

\indent Bien que ces technologies aient principalement été développées pour des environnements dits « structurés » (réseau routier et autoroutier), de nombreux autres secteurs expriment le besoin d'automatiser ou d'assister la conduite en milieu non structuré. Ces environnements complexes, souvent évolutifs et hors de portée des réseaux de communication conventionnels, présentent des défis scientifiques et techniques spécifiques. On peut citer l'agriculture de précision pour l'automatisation des exploitations, ou encore les opérations en milieux hostiles (chantiers, interventions sur incendies, opérations militaires), où une assistance décisionnelle avancée permet de décharger l'opérateur de certaines charges physiques et cognitives. C'est dans ce contexte que l'Agence de l'Innovation de Défense (AID), l'Agence Nationale de la Recherche (ANR) et la Direction Générale de l'Armement (DGA), en partenariat avec l'Agence d'Innovation pour les Transports (AIT) et le Centre National d'Études Spatiales (CNES), ont initié le Challenge MOBILEX. Étalée de fin 2023 à fin 2026, cette compétition oppose plusieurs consortiums académiques et industriels au travers d'épreuves à difficulté croissante, avec pour objectif d'accélérer la recherche et la montée en maturité technologique sur la thématique de la mobilité autonome des véhicules en environnement complexe*. 
C'est pour répondre à ce challenge que la société de Services et d'Ingénierie en Systèmes (SSIS) partenaire majeure du secteur automobile depuis les années 2000, SHERPA Engineering s'est associée à l'INRAE et à l'Institut Pascal au sein du consortium MOBITER. Cette synergie offre aux laboratoires l'opportunité d'évaluer leurs travaux de recherche à l'état de l'art sur une plateforme réelle, tandis qu'elle permet à SHERPA Engineering de développer et d'intégrer de nouvelles briques logicielles propriétaires, financées conjointement par l'ANR et par des fonds d'investissement en R&D de l'entreprise.\\

\indent Le présent rapport s'articulera autour de plusieurs axes. Nous débuterons par une présentation détaillée de SHERPA Engineering, suivie d'une description de la répartition des responsabilités au sein du consortium MOBITER. Nous examinerons ensuite les objectifs et les exigences du Challenge MOBILEX, ainsi que le matériel embarqué sur le robot autonome. La section suivante sera consacrée aux briques logicielles développées pour répondre aux attentes de la DGA. Enfin, nous presenterons les travaux réalisés au cours de notre stage, incluant la mise en place et réalisation de tests terrain, le développement d'un outil automatisé de génération de rapports de test terrain et l'intégration d'un nouveau planificateur local dans le module de suivi de lisière.
